% Cote 119: offprints, reports, courses, CV and letters (typed and printed part): archive cote 119, site page 269 (PDF page 270).
% typescript. Draft compared: none.
% AI-transcribed by Claude Opus 5.5 (claude-opus-5-5, adaptive thinking, effort high), Message Batches API, batch r119_full (msgbatch_01DwZNm5Y7YsKBWyAN4sFs8w), 2026-11-13; not yet checked by eye.
\scanpage{119/269}
\mspage{16}

pour ``thésauriser'' les meilleurs chercheurs. A cette condition serait favorisé
également le processus de ``décentralisation des compétences'', dont il a été
question depuis longtemps et que chacun appelle de ses voeux, voeux plus ou
moins sincères, il faut bien l'avouer, suivant que celui qui les fait se trouve
en province ou à Paris. Je serais certes le dernier à sous-estimer la force
d'une telle tentation pour un mathématicien passionné par sa science. Aussi
j'estime que la seule façon efficace d'y remédier serait la création d'une
\uline{Commission de mathématiciens interdépartements}, dont un des rôle serait de
superviser les nominations, à laquelle les départements de Mathématique
adhérents délégueraient une autorité véritable, y compris dans certains cas
à préciser un droit de veto, dans le but de garantir une répartition raison-%
nable des chercheurs suivant les diverses Universités, en tenant compte à la
fois de l'intérêt des divers centres de recherches effectivement existants dans
les divers départements, et des désirs des chercheurs concernés. Il s'agit
manifestement là d'un point extrêmement délicat, qui demandera des mécanismes
d'autocontrôle d'une grande souplesse (excluant bien entendu l'intervention
de toute instance ministérielle). Elles ne pourront évidemment fonctionner
\hfill que moyennant une évolution considérable
dans l'état d'esprit des mathématiciens concernés, et plus particulièrement
des parisiens, habitués à une concentration des compétences qui depuis long-%
temps dépasse de très loin les besoins particuliers de chacun de nous, et qui
à l'échelle nationale constitue une sorte de ``gaspillage par l'accumulation''.

8. Pour terminer, il convient de rappeler que jusqu'à aujourd'hui, l'autorité
à l'intérieur du département a été détenue exclusivement par les maîtres de
conférences et les professeurs, recrutés suivant les critères traditionnels,
c'est-à-dire grosso-modo, par des chercheurs. Ces derniers constituent une
véritable ``aristocratie éclairée'', et exercent un gouvernement absolu sur
une armée d'assistants et maîtres-assistants, aux effectifs quatre ou cinq
fois plus nombreux, et sur un ``sous-prolétariat'' d'étudiants. L'accession
d'un maître-assistant à l'aristocratie régnante est soumise à des critères
de recherche d'une rigueur telle que la grande majorité doivent y renoncer,
et par là même renoncer à un épanouissement normal de leurs capacités pédago-%
giques, comme membres responsables à part entière de la communauté des ensei-%

\margin{269}
